\documentclass{report}
\usepackage{amsmath,amssymb}
\setlength{\parindent}{0mm}
\setlength{\parskip}{1em}
\begin{document}
\begin{center}
\rule{6in}{1pt} \
{\large Austin Benson \\
{\bf Higher-order organization of complex networks}}

Stanford University \\ Gates Computer Science Building \\ Room 448 \\ 353 Serra Mall \\ Stanford \\ CA 94305 U S A
\\
{\tt arbenson@stanford.edu}\\
David Gleich\\
Jure Leskovec\end{center}

Two of the fundamental analyses in network science are clustering
(partitioning, community detection) and the frequency of network motifs,
or patterns of links between nodes. In this work, we develop spectral
algorithms to find network community structure defined by motifs. The
central computational kernel of our method is a large sparse eigenvalue
problem. We show several examples of motif-based communities in networks
from a variety of scientific domains.


\end{document}
